\sgaStableTarget{sga6.e05.prop.3-11}\textbf{Proposition 3.11.} \uline{Let} @@M0@@ \uline{be a} @@M1@@-ring, \uline{and} @@M2@@, @@M3@@ \uline{elements of} @@M4@@ \uline{such that} @@M5@@ \uline{for} @@M6@@. \uline{Put} @@M7@@; \uline{then}:
\begin{enumerate}[label=\roman*)]
\item \uline{for} @@M8@@, we have @@M9@@
\item \uline{for} @@M10@@, @@M11@@ \uline{is the} @@M12@@-th \uline{symmetric function of the} @@M13@@.
\end{enumerate}
Let @@M14@@. We then have @@M15@@. Since @@M16@@ for @@M17@@, and @@M18@@ for @@M19@@, @@M20@@. The proposition then follows from @@M21@@.

\sgaStableTarget{sga6.e05.sec.4}\section*{4. @@M22@@-rings generated by generators and relations}

We still work over @@M23@@.

\sgaStableTarget{sga6.e05.para.4-1}4.1. Let @@M24@@ be a @@M25@@-ring. Consider a family of indeterminates @@M26@@, and let @@M27@@; we wish to define on @@M28@@ a @@M29@@-ring structure extending that of @@M30@@.

To specify a homomorphism @@M31@@, it suffices to specify the images of the @@M32@@, @@M33@@, which defines a ring homomorphism: @@M34@@; for it to define a @@M35@@-ring structure on @@M36@@, we must verify the commutativity of diagram (3.7.1). We therefore put:
\sgaStableTarget{sga6.e05.eq.4-1-1}@@M37@@
Since diagram (3.7.1) is a diagram of ring homomorphisms, it suffices to verify its commutativity on the @@M38@@, @@M39@@ (it is known by hypothesis on the elements of @@M40@@); this follows readily from definitions (4.1.1) and the fact that @@M41@@ and @@M42@@ are @@M43@@-homomorphisms.

